\documentclass{article}


\begin{document}


\subsubsection{Flexible use}
\label{(1) Flexible use}

Stations can suffer from fluctuations of use, including overcrowding during peak times and emptiness during non-peak times (See \ref{Introduction} \hyperref[Introduction]{Introduction}).

% During peak times or when there are failures of railway lines, station areas can be overcrowded with detained passengers. 

\textbf{Solutions}: Flexible use is the solution that addresses fluctuation in use (fig. \ref{fig:flexible_use}). Two types of spaces can be flexibly used: 1) Stations' own spaces, and 2) City neighborhoods' spaces near the stations. The ideal scenarios are: During stations' peak times or emergent times, redundant spaces are used for accommodating passengers; during non-peak times, redundant spaces are used for city leisure activities or events (such as open markets, outdoor music festivals, and carnivals).

Stations need to \hyperref[Connect with neighborhoods]{Connect with neighborhoods (8)}, so that spaces of both cities' spaces can be flexibly used by stations accommodating peak-time passengers. For non-peak-time leisure use, it is necessary to further \hyperref[(22 23) Reduce barriers to ease flow]{Reduce barriers to ease flow (22) \& (23)}, provide \hyperref[(7) Human-oriented spaces]{Human-oriented spaces}, with leisure events and \hyperref[(6) Make events visible by aligning main paths and open spaces]{Make events visible by aligning main paths and open spaces (6)}.

Station areas need to also to be \hyperref[(12) Programming considering the time dimension]{Programmed considering the time dimension (12)} for flexible use of spaces.

Stations can be designed as \hyperref[(14 15) Flexible buildings or building components]{Flexible buildings or building components (14) \& (15)}, to have a responsive capacity for accommodating the peak time passengers. 

In some cases, railway stations are used temporarily, or temporarily drastically increase capacity, to facilitate city events or nearby stadiums, such as the Rotterdam Stadion Station (for Feyenoord Stadium), Victoria Station Manchester (for AO Arena), Zandvoort Station (for Formula 1 car racing events), Bijlmer Arena (for multiple nearby venues including Johan Cruyff Arena and Ziggo Dome). However, this flexible change of capacity is more of a pure managerial intervention rather than a spatial design intervention.

\begin{figure}[h]
    \centering
    \includegraphics[width=0.5\columnwidth]{Paper3Imgs_pdfs/index_pictures.027.pdf}
    \caption{Flexible use}
    \label{fig:flexible_use}
\end{figure} \FloatBarrier

\textbf{Concepts clarifications}: Of notice, in the 'flexible use,' the 'time dimension' is intrinsically considered, while some similar terms do not necessarily incorporate it. To better emphasize the \textbf{time dimension} nature, terms like 'flexible,' 'temporary,' 'reconfigurable,' 'adaptive,' and 'changeable' are more appropriate than terms like 'mixed-use,' 'multi-functional,' 'all-in-one,' and 'sharing.' Because the latter can mean something has multiple functions working at the same time, not necessarily at different times. In the practice of urban planning and metro station development, similar to 'flexible use,' the concept 'mixed-use' \citep{rowleyMixeduseDevelopmentAmbiguous1996} has long been promoted. But 'mixed-use' is not evident to readers as relates to time, even though a few scholars also incorporate the time dimension \citet{hoppenbrouwerMixeduseDevelopmentTheory2005} as the fourth dimension, beyond the 3D physical spatial dimensions, into the 'mixed-used' definition. 

Flexible sharing with neighborhoods at least has two positive implications: 1. Given the same investment, station-city will have larger capacities. 2. Given the same capacities, smaller spaces are needed, which means less investment and less redundant space (and hence less emptiness during non-peak times).

\textbf{Limitations}: Flexible use has its limitations, as it can potentially add to the management efforts for stakeholders of stations and cities (e.g., see the chaos in Bijlmer Arena station when the management capacity is surpassed \citep{aseniyaCrowdControlBijlmer2019}), and as it can exacerbate overcrowding and emptiness if time dimension is not well considered.








\subsubsection{(2) Suitable general layouts of the station and city}
\label{(2) Suitable general layouts of the station and city} 

\textbf{Context \& Problems}: Stations and cities have different elements (and can be seen as the mass and the void); depending on the contexts, a suitable general layout of them can make the spatial allocation more understandable and rational. The following are some options of the general layouts: 
% - overcrowding
% - emptiness

\textbf{Solutions \& Examples}: The general layouts of stations' and cities' elements can be 
1) separated;
2) connected - e.g., Rotterdam Central Staiton and Utrecht Central Staiton, where stations are \hyperref[Connect with neighborhoods]{Connect with neighborhoods (8)} through \hyperref[(25) City axes]{City axes (25)}, and \hyperref[(20 21) City passages and their positioning]{(20) City passages and their positioning};
3) stacked - e.g., Montparnasse Station, upon which are apartment and offices buildings, and rooftop parks;
4) mixed - e.g., Shibuya Station and Shinjuku Station, where stations' spaces and cities' spaces are inseparable;
(fig. \ref{fig:generallayout}) \citep{shikataIntegratedStationCityDevelopment2013, qiAnalysisStationCitySynergetic2019, dewildeMultipleUseSpace2006}.

Whether mixing or separating stations' and cities' elements depends on the specific project contexts. Bringing non-transport city functions into the station areas, has several benefits, including intensifying land use, more commercial revenue, vibrating the spaces, and so on. However, this can potentially also cause conflicting use between the stations and cities, between passengers and non-passengers. Therefore, sometimes it is necessary to \hyperref[(3) Set apart non-transport function]{Set apart non-transport function (3)}.

These general layouts can be applied on different spatial scales, as long as the layouts facilitate designers' or stakeholders' understanding of the spatial configurations of station areas/buildings/elements.

\textbf{Limitations}: There are some limitations of trying to apply the general layouts: 1) The layouts can be meaningless or merely abstract concepts if they not related to the actual functioning of the stations' cities' components. 2) Also, sometimes defining stations and city elements is hard work itself. These elements may not have clear identifiable boundaries, or they have \hyperref[(1) Flexible use]{Flexible use (1)}.

\begin{figure}[h]
    \centering
    \includegraphics[width=0.5\columnwidth]{Paper3Imgs_pdfs/index_pictures.020.pdf}
    \caption{Different general layouts of the station and city}
    \label{fig:generallayout}
\end{figure} \FloatBarrier







\subsubsection{(3) Set apart non-transport function}
\label{(3) Set apart non-transport function}

\textbf{Contexts}: Setting apart non-transport function is one among several \hyperref[(2) Suitable general layouts of the station and city]{Suitable general layouts of the station and city (2)}. 

In the station-city integration practice, and the related and more popular topic of transit oriented development (TOD), mixed use is a fundamental principle. Despite the the mixed-use strategy have many benefits, in some situations, it also causes problems, and these problems can be easily overlooked.

\textbf{Problems}:
Mixing of stations with cities, if not done properly, can make the station buildings unrecognizable, and can bring interference in flow between passengers and non-passengers.

% in Section \ref{Spatial scales of concern: the district and building levels} \hyperref[Spatial scales of concern: the district and building levels]{Spatial scales of concern: the district and building levels}, 


\textbf{Solutions \& Example}:
Therefore, setting apart non-transport function is a solution in the above situation. One representative example is the Utrecht Central Station area. Before the latest redevelopment, the Utrecht Central Station building is highly mixed with the adjacent Hoog Catharijne shopping mall building. This mixture raised complaints from passengers as the station building is unrecognizable, passengers felt lost once they got off the train arriving this area, and the passages across the station building was always overcrowded with passengers and city users. During the latest development, a east well-designed plaza is added, separating the station building from the shopping mall building (fig. \ref{fig:utrechtsetapartnonstation}), and endowing the station a clear image identity. A spacious footbridge, with street-like qualities, is built alongside the station building for pedestrians to pass-by freely.

\textbf{Limitations}: The distinction between separation and mixing can be ambiguous, so apply this design principle (i.e., Setting apart non-transport function) still depends much on designers' experience.

\begin{figure}[h]
    \centering
    \includegraphics[width=0.5\columnwidth]{Paper3Imgs_pdfs/index_pictures.016.pdf}
    \caption{Set apart non-transport function}
    \label{fig:utrechtsetapartnonstation}
\end{figure} \FloatBarrier




\subsubsection{(4) \& (5) Vibrate city environment or increase accommodating capacity by scattered mobility nodes}
\label{(P4 P5) Vibrate city environment or increase accommodating capacity by scattered mobility nodes}

\textbf{Contexts \& Problems \& Solutions}: A layout of scattered mobility nodes can be applied for the following purpose:
1) To let railway passengers, who are potential customers of city stores, go through city districts, vibrating the city environment.
2) To increase the peak time capacity of the mobility sites -  passengers are more evenly dispersed to multiple sites - instead of concentrating on a limited site and, therefore, are less prone to overcrowding or even stampedes.
% increase capacity (Building-level, P5) (fig. ). Given the same volume of pedestrians, a more scattered layout of mobility sites enables less dense concentrations of pedestrians in each site. 

Scattered layouts of mobility nodes can be applied at the district level or the building level (figs. \ref{fig:districtscattered}, \ref{fig:buildingscattered}). 

\textbf{Examples}: Cases like Shanghai railway station adopt the layout of scattered mobility sites. 

\textbf{Limitations}: The side effects include: 1) that more mobility sites result in harder wayfinding, and 2) harder transferring between mobility sites (Therefore, if a station has a heavy transfer load between different mobilities, it is fine to use a scattered layout of mobility nodes; While if it has light transfer load between different mobilities, it is not appropriate to use this layout).

\begin{figure}[h]
    \centering
    \includegraphics[width=0.5\columnwidth]{Paper3Imgs_pdfs/index_pictures.017.pdf}
    \caption{Vibrate city environment by scattered mobility nodes}
    \label{fig:districtscattered}
\end{figure} \FloatBarrier
\begin{figure}[h]
    \centering
    \includegraphics[width=0.5\columnwidth]{Paper3Imgs_pdfs/index_pictures.014.pdf}
    \caption{Increase transport capacity by scattered mobility nodes}
    \label{fig:buildingscattered}
\end{figure} \FloatBarrier




\subsubsection{(6) Make events visible by aligning main paths and open spaces}
\label{(6) Make events visible by aligning main paths and open spaces}

\textbf{Contexts}: If a station suffers from emptiness, to vibrant the environment and reduce emptiness, leisure events can be organized.

\textbf{Problem}: But not every space is supportive of leisure events. If events happen in some places that are invisible to pedestrians (e.g., due to level changes, street-level sight lines are blocked) or away from the main path, it will be hard for events to attract enough people, and therefore, the empty station spaces remain empty. For instance, in fig. \ref{fig:pathwithevent}, if events are held in both stations, the environment in the right scenario will be energized significantly while the environment in the left won't.

\textbf{Solutions \& Examples}: The alignment between open spaces and event-goers' (leisure users') paths should be well-designed, to make sure events are visible, or even accessible for pedestrians to participate - depending on the event types. In the design stage, designers can either align spaces with paths or vice versa, aligning pedestrians' paths with spaces. 

Many different types of events can be held in station areas. For example, open markets, outdoor music festivals, carnivals, and so on \citep{smithStagingCityEvents2021}
\textbf{Activities} can be seen as small-scale events. 

The events' visibility can be enhanced by using \hyperref[(24) Stairs as stages or seats]{Stairs as stages or seats (24)}.

The event-supportive spaces are usually \hyperref[(7) Human-oriented spaces]{Human-oriented spaces (7)} that support \hyperref[(1) Flexible use]{Flexible use (1)}.

The paths that event-goers (leisure users) are willing to walk through should be considered in a large network of paths, which can consist of \hyperref[Connect with neighborhoods]{Connections with neighborhoods (8)} that \hyperref[(22 23) Reduce barriers to ease flow]{(22 23) Reduced barriers to ease flow}, and \hyperref[(20 21) City passages and their positioning]{City passages and their positioning (20) \& (21)}.

% \textbf{Limitations}: 

\begin{figure}[h]
    \centering
    \includegraphics[width=0.5\columnwidth]{Paper3Imgs_pdfs/index_pictures.001.pdf}
    \caption{Alignment between open spaces and main paths}
    \label{fig:pathwithevent}
\end{figure} \FloatBarrier







\subsubsection{(7) Human-oriented spaces}
\label{(7) Human-oriented spaces}

% - flexible use - non-transport activities; leisure events

\textbf{Contexts}: It is debatable if station areas should prioritize human uses, as stations are fundamental transportation infrastructures, and hence, people would think of stations as primarily for vehicles. This debate exists in the urban planning field for the general urban environment and for other transport infrastructures, including streets \citep{bertoliniStreetsTrafficStreets2020}.

In the past decades, there has been a paradigm shift from car-oriented to human-oriented planning; and for the design of railway station areas, at least in several situations, human-centered spaces should be a priority: 1) When station areas are shared with the cities and used by non-transport users (see principle \hyperref[(1) Flexible use]{Flexible use (1)}), or 2) when it is necessary to organize leisure events to reduce space emptiness (see chapter \ref{Introduction} \hyperref[Introduction]{Introduction}).

\textbf{Problems}: To support city non-transport use or leisure events, station spaces should be attractive. However, the human experience is often neglected in the design of some stations, where problematic features make them unattractive to people: 1) traffic elements dominating the spaces in railway station areas; 2) open spaces surrounded by traffic roads; 3) station buildings and spaces being monolithic, which makes them visually boring, and hard for people to move on foot for the long distance between facilities.

\textbf{Solutions \& Examples}: Human-oriented spaces should have good spatial qualities that enhance human experience, as many as possible. The qualities depend on different cultural contexts, project contexts, different users with different evolving demands, and so on. It is impossible to exhaustively define all the qualities, and that is perhaps why Alexander calls it the 'quality without a name' \citep{alexanderTimelessWayBuilding1979}. Nevertheless, in urban design field, scholars have explicated some common qualities or spatial features, for instances: \textbf{human scale}; \textbf{enclosed} spaces with well-defined boundaries \citep{sitteCityPlanningAccording1889}; positives \textbf{edges} that promote social activities, urban life \citep{gehlLifeBuildings1987, whyteSocialLifeSmall2001}, and recognition/identity of different areas \citep{jacobsDeathLifeGreat1961, lynchImageCity1960}; \textbf{complexity/diversity} of visual elements \citep{ewingMeasuringUnmeasurableUrban2009}. 

The spatial layout, besides space qualities, also matters for human experience.  In fig. \ref{fig:humancenteredspace}-left, we try to illustrate some features regarding the layout: a human-centered space should be a coherent space at a core position in the station areas or the station building, free from traffic elements, and free from intensive pass-by transport flow so that people are possible to stay there. In contrast, the layout in fig. \ref{fig:humancenteredspace}-right does not have these features.

A significant example is the station plaza of the new Rotterdam Central Station, as against the old one (fig. \ref{fig:Rot_old_new}). In the old scenario, the core area was occupied by transport elements, including trams, buses, car parking, and a metro station entrance building; Walkable spaces are located around this core area; Hardly any area is possible for people to stay. In the new scenario, the core area is designed for humans; Metro entrances, tram platforms, and bus stations are moved aside; Car parking is located underground. During large events such as the Rotterdam Summer Carnival, the new station plaza is used without inference to the normal functioning of the station's transport operations. While in the old scenario, assuming there would have been an event, the old plaza was hardly possible to use unless all the transport elements stopped functioning.

Addressing high fluctuations of use towards demand-responsive station areas, human-oriented spaces can be supplemented with several smaller design principles, including \hyperref[(9 10) Smooth level changes by landscape design]{Smooth level changes by landscape design (9)}, \hyperref[(11) Adaptive redundant spaces with path regulation]{Redundant spaces with path regulation (11)}, \hyperref[(16) Add installations and facilities]{Add installations and facilities (16)} and \hyperref[(17 26) Reconfigurable elements or spaces]{Reconfigurable elements or spaces (17) \& (26)}.

Pedestrian paths, although usually linear elements for passing through, can be designed as human-oriented spaces if these paths are also supposed to let people stay. From this perspective, human-oriented spaces can be used as a strategy to support \hyperref[(25) City axes]{City axes (25)} and \hyperref[(22 23) Reduce barriers to ease flow]{Reduce barriers to ease flow (22) \& (23)}.

\textbf{Limitations}: There are some limitations of human-oriented spaces, including 1) the possible higher cost to create such spaces; 2) certain good qualities in one case may not necessarily be good in another case.

\begin{figure}[h]
    \centering
    \includegraphics[width=0.5\columnwidth]{Paper3Imgs_pdfs/index_pictures.002.pdf}
    \caption{Human-oriented spaces versus vehicle-oriented spaces}
    \label{fig:humancenteredspace}
\end{figure} \FloatBarrier
\begin{figure}[h]
    \centering
    \includegraphics[width=0.5\columnwidth]{Paper3Imgs_pdfs/index_pictures.029.pdf}
    \caption{The old and the new Rotterdam central stations}
    \label{fig:Rot_old_new}
\end{figure} \FloatBarrier



\subsubsection{(8) Connect with neighborhoods}
\label{Connect with neighborhoods}

\textbf{Contexts}: In the following context, stations should connect with their city neighborhoods:
% - lack of space, sharing to increase the capacity during peak times 
% - flexible use - promote movement, energize the environment
% - when city axes are already planned

% 1) If a station has high fluctuations of use, it should connect with its neighborhoods so that: During peak times or emergent, the station can use city spaces for high capacity accommodating strained passengers. During non-peak times, city users can use the redundant spaces in the stations. 

When stations have high fluctuation of use, and \hyperref[(1) Flexible use]{Flexible use (1)} with neighborhoods is needed, for the purpose of temporarily using the neighborhoods' spaces to accommodating stations' strained passengers during peak or emergent times; 

\textbf{Problems}: In contrast, if there are no such connections, problems will occur: the station spaces suffer from overcrowding during peak or emergent times.


\textbf{Solutions}: Connecting with the neighborhoods with stations (fig. \ref{fig:connection}), therefore, becomes a solution. Connections/linkages include pedestrian paths, footbridges/skywalks, tunnels, and so on. Passengers can use these connections when they have to. 

Connections are typically smaller parts of longer paths. They often cause sudden changes in the width of the paths, so it is necessary to ensure that they are not bottlenecks of flow. The connections should have enough capacity for passengers to pass during peak times.

The connections can basically serve for pass-through usage, whereas more preferably, the connections should be designed in a pleasant way that \hyperref[(22 23) Reduce barriers to ease flow]{Reduce barriers to ease flow (P22)} (Unless the stations do NOT want attractive non-transport city users).

% Causes of peak flow: the surge of passengers due to events, holidays, festivals …Bottlenecks should be reduced as much as possible.

Within the full pedestrian paths in which stations and city neighborhoods are connected, \hyperref[(20 21) City passages and their positioning]{City passages and their positioning (20) \& (21)} should also be considered.


% \textbf{Examples}:
% Rotterdam, Utrecht, negative examples: Beijing West, 
% \textbf{Limitations}:


\begin{figure}[h]
    \centering
    \includegraphics[width=0.5\linewidth]{Paper3Imgs_pdfs/index_pictures.010.pdf}
    \caption{Connect with neighborhoods}
    \label{fig:connection}
\end{figure} \FloatBarrier




\subsubsection{(9) \& (10) Smooth level changes by landscape design}
\label{(9 10) Smooth level changes by landscape design}

% - bottleneck adds to overcrowding
% - barrier of leisure flow - level changes that draw too much attention of pedestrians


\textbf{Context \& Problems}: Level changes are common in station areas due to many station buildings being multi-floor, or having topographic variations in the surrounding landscape. Using stairs naturally becomes a solution for level changes (see also design principle \hyperref[(24) Stairs as stages or seats]{Stairs as stages or seats (24)} can also be considered). However, walking on stairs is not a smooth experience - it requires careful attention from pedestrians in case of falling down, and it is even more dangerous during overcrowding. To further advance \hyperref[(7) Human-oriented spaces]{Human-oriented spaces (7)}, alternatives are needed.

\textbf{Solutions \& Examples}:
Using landscape design (figs. \ref{fig:jiaxing}, \ref{fig:arnhem}) helps to smooth the level changes, as compared to stairs. It would be useful, especially in flow-intensive areas, to prevent pedestrians from falling down. Landscape design can apply on both the district level (e.g., see fig. \ref{fig:jiaxing}, the plaza in the early design proposal of Jiaxing Railway Station, China, by MAD Architects) and the building level (e.g., the slop in the station hall of Rotterdam Central Station). When applied at the building level, it can even bring fresh spatial experience (e.g., Arnhem Central Station fig. \ref{fig:arnhem}).

\textbf{Limitations}: Landscape design smoothing level change can be expensive, especially at the building level. (The design of buildings like Arnhem Central Station is usually related to the so-called 'parametric design,' which requires computer algorithms dealing with complex geometries with many engineering challenges). Simply using escalators could be a cheaper alternative.

landscape design smoothing level change may cause confusion in wayfinding for people who are not familiar with such level-changing environments.

\begin{figure}[h]
    \centering
    \includegraphics[width=0.5\columnwidth]{Paper3Imgs_pdfs/index_pictures.007.pdf}
    \caption{Smooth level changes by landscape design at the district level}
    \label{fig:jiaxing}
\end{figure} \FloatBarrier
\begin{figure}[h]
    \centering
    \includegraphics[width=0.5\columnwidth]{Paper3Imgs_pdfs/index_pictures.022.pdf}
    \caption{Smooth level changes by landscape design at the building level}
    \label{fig:arnhem}
\end{figure} \FloatBarrier





\subsubsection{(11) Adaptive redundant spaces with path regulation}
\label{(11) Adaptive redundant spaces with path regulation}

\textbf{Contexts \& Problems}:
Adaptive redundant spaces with path regulation should be used if a station has the following dilemma:
1) The station has \hyperref[(18 19) Redundant spaces or setting apart bottlenecks]{Redundant spaces or setting apart bottlenecks (18 19)} to reduce overcrowding during peak times; But
2) during non-peak times, the added long transfer distance by redundant spaces annoys passengers.

% Redundant spaces can make transfer distance longer if they are directly located between originations and destinations. So it seems a dilemma to retain redundant space for peak time use and ensure efficient transfer during non-peak time; 

\textbf{Solutions \& Examples}: Regulating paths solves the above peak-time redundancy and non-peak-time transfer efficiency dilemma. It is similar to \hyperref[(P13) Shortcuts or optimizing paths]{Shortcuts or optimizing paths (P13)}.

During peak times, it is a common practice to regulate the path by fences or lines (fig. \ref{fig:fencezigzag}) to ensure pedestrians move in an orderly queue, which prevents pedestrians from concentrating randomly, leading to stampedes. 
% Referring to this practice, during peak time, using fence-like measures can also regulate the path to make people walk through the redundant spaces, where they do not go during non-peak time, to avoid overcrowding. 

Instead of fences or lines, which are quite intrusive elements, paths can also be regulated by more user-friendly, even interactive design elements, such as urban furniture, \hyperref[(17 26) Reconfigurable elements or spaces]{Reconfigurable elements (17)} - water bodies, fountains, potted plants, art installations, stalls, and so on (fig. \ref{fig:Path_regula_reconfig_elements}).

To further satisfy the waiting passengers in the regulated zigzag paths, it is beneficial to create a pleasant environment of staying, compensating for the longer waiting time than usual.

During non-peak times, the path-regulating elements can be removed to let pedestrians move freely. \hyperref[(7) Human-oriented spaces]{Human-oriented spaces (7)} is a prerequisite for redundant space to be adaptively used during non-peak times.

In Chinese station cases, path regulations (using fences) are more commonly seen during peak times. In the Netherlands, however, it is less common. 

\textbf{Limitations}: Path regulations can easily raise complaints from passengers or nearby citizens; therefore, they should be used sparingly with strong necessity. 

\begin{figure}[h]
    \centering
    \includegraphics[width=0.5\columnwidth]{Paper3Imgs_pdfs/index_pictures.012.pdf}
    \caption{Adaptive redundant spaces with path regulation}
    \label{fig:fencezigzag}
\end{figure} \FloatBarrier
\begin{figure}[h]
    \centering
    \includegraphics[width=0.5\columnwidth]{Paper3Imgs_pdfs/index_pictures.019.pdf}
    \caption{Path regulation using reconfigurable elements}
    \label{fig:Path_regula_reconfig_elements}
\end{figure} \FloatBarrier





\subsubsection{(12) Programming considering the time dimension}
\label{(12) Programming considering the time dimension}

\textbf{Contexts \& Problems}: To reduce overcrowding and emptiness in stations, \hyperref[(1) Flexible use]{Flexible use (1)} is needed, and one planning strategy is programming station areas considering the time dimension (fig. \ref{fig:programmingsharing}).

\textbf{Solutions \& Examples}: Green parks and pedestrian-friendly streets, if programmed next to stations, can be temporarily restrictedly used by stations exclusively - for accommodating waiting passengers during railway peak times or for evacuation during emergent times. 

Similarly, large infrastructures such as sports stadiums, and music venues, since they are usually not used on a daily basis, can be shared with stations for railway peak time use. For example, in Guangzhou, China, multiple sports stadiums were planned for backup use of railway stations if any emergency occurred during Chinese New Year times \citep{xinhuanewsagencyGuangzhouSpringFestival2016}. This makes sense in this culture, as during the Chinese New Year, there are surges of railway passengers \citep{mccarthyChineseNewYear2018} who want to go to their hometowns for family gatherings, while at the same time, sports stadiums are mostly empty.

During non-peak times, the city spaces and infrastructures can be used as normal for their original non-transport leisure activities. Or, they can even take advantage of the railway as a more sustainable means of transportation \citep{newmanSustainabilityCitiesOvercoming1999}. For example, in Amsterdam, the Netherlands, several stadiums are planned next to Bijlmer Arena railway station, promoting the use of the railway. 


\begin{figure}[h]
    \centering
    \includegraphics[width=0.5\columnwidth]{Paper3Imgs_pdfs/index_pictures.018.pdf}
    \caption{Programming considering the time dimension (for flexible use)}
    \label{fig:programmingsharing}
\end{figure} \FloatBarrier

\textbf{Limitations}: Programming has its limitations for events with frequency from multi-weekly to yearly (which are our focuses in this research). We discovered only limited programming strategies for this kind of event. In contrast, for more frequent events such as commuting (with a frequency between daily to weekly, and is an issue more for metro stations instead of railway stations), there are common programming strategies that have long been established, such as mixing housing/hospitality with offices.






\subsubsection{(13) Shortcuts or optimizing paths}
\label{(P13) Shortcuts or optimizing paths}

% (This design principle is more suitable for large Chinese stations)
\textbf{Context \& Problems}: In some stations, the transfer between different mobilities is effect-consuming. For instance, at Beijing West Railway Station, if a passenger comes out of the metro and wants to go to the railway, he needs to check the ticket out of the metro, go the lengthy paths with level changes, and then check in at the railway with a strict security check, and a ticket check (fig. \ref{fig:buildshortcut}). Therefore, shortcuts in the paths or optimizing paths are needed. 

\textbf{Solutions \& Examples}: If there are shortcuts (direct links between the metro and the railway) in the paths (fig. \ref{fig:buildshortcut}), the passenger can save transfer distance and be relieved from double security checks (He already had a security check when he entered the metro system). For example, Beijing South Railway Station (RS), as a counterpart of Beijing West RS, provides direct transfer links between the metro and the railway. 

% If there is a shortcut enabling him to directly go from the metro station to the railway station (e.g., refer to Beijing South Station), his transfer experience (in terms of transfer time and distance) will be improved, and there will be less redundant space (hence overcrowding during peak times and emptiness during non-peak times) needed for passengers like him.

The shortcuts can be applied during non-peak times for transfer speed, while shut down during peak times for safety (as a lengthy path has more redundant spaces to accommodate the strained passengers, preventing overcrowding).

Shortcuts in paths are similar to \hyperref[(11) Adaptive redundant spaces with path regulation]{Path regulation for adaptive redundant spaces}. Their difference (in this research) is that shortcuts are more for (indoor) buildings, and path regulation is more for outdoor spaces (fig. \ref{fig:Pattern_language_network}).

\textbf{Limitations}: However, adding shortcuts means adding extra ingress/egress paths, in which extra electronic systems are needed, hence increasing costs for railway companies. It also adds more complexity for the system and for the users. The well-designed \hyperref[(7) Human-oriented spaces]{Human-oriented spaces (7)} along the normal paths will be missed if passengers choose the shortcuts over the normal paths.

\begin{figure}[h]
    \centering
    \includegraphics[width=0.5\columnwidth]{Paper3Imgs_pdfs/index_pictures.006.pdf}
    \caption{Shortcuts or optimizing paths}
    \label{fig:buildshortcut}
\end{figure} \FloatBarrier






\subsubsection{(14) \& (15) Flexible buildings or building components}
\label{(14 15) Flexible buildings or building components}

\textbf{Contexts \& Problems \& Solutions}:
If a station has huge fluctuations of passengers, to support \hyperref[(1) Flexible use]{Flexible use (1)}, they can be designed as a flexible building, or a building with some flexible components (including, such as the flexible gate machines). 

Similar to the word 'flexible,' temporarily used buildings solve the emptiness issue from the root. There are various types of buildings titled with concepts that enable a certain degree of flexibility. These concepts include 'changeable,' 'movable,' or 'temporary buildings,' 'flexible,' 'multi-functional' (in temporal dimension), 'reconfigurable,' 'kinetic' \citep{parametricarchitectureWhatKineticArchitecture2024}, and so on \citep{geraedtsSuccessFailureFlexible2011}. These buildings have flexibility either by changing their use (programming) during different times, or their existence (assemblage) during certain periods. 

\textbf{Examples}:
Flexible buildings in the real world are not common due to the high cost and high maintenance. Here are a few non-station examples: 1) The Shed (fig. \ref{fig:theshed}), which has a movable hall as an event venue. 2) Stadium 974 (fig. \ref{fig:Changeable_building}), which is a stadium made of containers and can be assembled and shipped (see also design proposal by \citet{smithShipStadiumTransforming2023}). For station cases, there are examples such as 3) temporary station buildings during the construction of the main new station buildings (e.g., in the Rotterdam Central Station case). In Shanghai South Railway Station and Beijing West Stations, many redundant checking machines are allocated in the gate areas, which can be used during peak times for higher capacity.

\textbf{Limitations}: Flexible buildings are expensive to build and maintain; They are like a fantasy against reality in most situations.

\begin{figure}[h]
    \centering
    \includegraphics[width=0.5\columnwidth]{Paper3Imgs_pdfs/index_pictures.008.pdf}
    \caption{A building with changeable building component - the Shed}
    \label{fig:theshed}
\end{figure} \FloatBarrier
\begin{figure}[h]
    \centering
    \includegraphics[width=0.5\columnwidth]{Paper3Imgs_pdfs/index_pictures.024.pdf}
    \caption{A changeable building - Stadium 974}
    \label{fig:Changeable_building}
\end{figure} \FloatBarrier




\subsubsection{(16) Add installations and facilities}
\label{(16) Add installations and facilities}

\textbf{Contexts \& Problems}: In the following contexts, it will be helpful to add installations and facilities to vibrant the environment: 
1) If a station suffers from emptiness during non-peak times; and 
2) the station spaces are huge and boring.
3) The station spaces are human-oriented spaces that are supportive of events leisure activities or big leisure events

\textbf{Solutions}:
Adding installations and facilities (fig. \ref{fig:antwerpferrywheel}) is helpful to reduce emptiness by 1) stimulating events; 2) reducing the size of huge empty spaces to make them more human-scale; 3) making the scenes visually more diversified to have enough complexity. Installations can be interactive (e.g., ferry wheels) or just visually interesting (big inflatable art installations \citep{iversStagingUrbanLandscapes2018}).

The installations and facilities can also be \hyperref[(17 26) Reconfigurable elements or spaces]{Reconfigurable elements (17)} that can be rearranged. 

\textbf{Examples}:
For example, At Antwerpen Central Station, a ferry wheel is sometimes installed in the front square. At Rotterdam Central Station, different art installations are installed from time to time, on the front square and in the station hall.

\textbf{Limitations}: If the installations and facilities are not attractive by themselves, then they can deteriorate the space quality.
For instance, in the Antwerpen Central Station case, when the ferry wheel is installed, the corresponding facilities, such as ticket-selling points, can make the square messy. On the street level, pedestrians' views are much blocked by these facilities, and the beautiful facade of the station is less visible to pedestrians. Similarly, on the South Square of Beijing West Railway Station, many staircases, soldier standing points, the roofs of sunken atriums, police car parking points, and so on, also block pedestrians' views.

\begin{figure}[h]
    \centering
    \includegraphics[width=0.5\columnwidth]{Paper3Imgs_pdfs/index_pictures.003.pdf}
    \caption{Add installations and facilities}
    \label{fig:antwerpferrywheel}
\end{figure} \FloatBarrier




\subsubsection{(17) \& (26) Reconfigurable elements or spaces}
\label{(17 26) Reconfigurable elements or spaces}

\textbf{Contexts \& Problems}: Reconfigurable elements or spaces can be used for the following purposes:
1) To vibrate the environment if station spaces are empty during non-peak times.
2) To \hyperref[(11) Adaptive redundant spaces with path regulation]{Regulate path (17)} during peak times.
3) To further develop \hyperref[(7) Human-oriented spaces]{Human-oriented spaces (7)}, making these spaces more diverse with elements.
4) To make \hyperref[(16) Add installations and facilities]{Installations and facilities (16)} more flexible

\textbf{Solutions \& Examples}:
Reconfigurable elements are movable. They can be rearranged in different layouts for different use scenarios. They can be urban furniture (fig. \ref{fig:urbanfurniture}), potted plants, art installations, stalls, and so on. Reconfiguration of elements can be a spontaneous activity of users, such as people who use urban furniture will move the furniture by themselves.

Reconfigurable spaces can be set as different venues for different activities or events
(fig. \ref{fig:Reconfigurable_spaces}). For example, Bryant Park in New York is constantly reconfigured as a music venue, tennis courts, a grassland, and an ice skating rink \citep{iversStagingUrbanLandscapes2018}. 

Reconfigurable spaces or elements are similar to \hyperref[(14 15) Flexible buildings or building components]{Flexible buildings or building components (14) \& (15)}. Their difference is that the former is about outdoor spaces, and the latter is about (indoor) buildings.

\begin{figure}[h]
    \centering
    \includegraphics[width=0.5\columnwidth]{Paper3Imgs_pdfs/index_pictures.004.pdf}
    \caption{Reconfigurable elements}
    \label{fig:urbanfurniture}
\end{figure} \FloatBarrier
\begin{figure}[h]
    \centering
    \includegraphics[width=0.5\columnwidth]{Paper3Imgs_pdfs/index_pictures.026.pdf}
    \caption{Reconfigurable spaces}
    \label{fig:Reconfigurable_spaces}
\end{figure} \FloatBarrier


\subsubsection{(18) \& (19) Redundant spaces or setting apart bottlenecks}
\label{(18 19) Redundant spaces or setting apart bottlenecks}

\textbf{Contexts \& Problems}: In the following situation:
Pedestrians move between different mobility sites in station areas, from originations to destinations, and there may be bottlenecks(Section \ref{Viewing stations from the pedestrian flow perspective}). 
1) If a station has surges of passengers during peak times; and 
2) if the mobility sites or bottlenecks are too close, then the strained passengers will quickly be added up to overcrowding.
Therefore, it is necessary to add redundant spaces or, equivalently, set apart bottlenecks.

\textbf{Solutions}: Redundant spaces (between originations to destinations, between different mobility sites, or between bottlenecks) reduce overcrowding in two ways:
1) Redundant spaces make pedestrians walk longer. Since different pedestrians walk at different speeds, the surges of pedestrians at the originations will be smoother when the pedestrians arrive at the destinations.
2) Redundant spaces are extra spaces that accommodate strained pedestrians.
Redundant spaces can be set at the district level (fig. \ref{fig:layout_setapart}) as well as at the building level (fig. \ref{fig:build_setapart}).

\textbf{Examples}: One example of redundant spaces being used is the Allianz Arena Munich. From the nearest train station to the stadium, it takes over 1km distance to walk. The massive space in between is redundant space that can accommodate a large number of people.
% Setting apart bottlenecks is similar to adding redundant space in between. If two bottlenecks are located too close, detained pedestrians caused by each of them are further concentrated, making the situation worse. Setting the two apart allows more (redundant) space and capacity to accommodate the aggregation of detained pedestrians.

\textbf{Limitations}: Redundant spaces have the following limitations:
1) Despite redundant spaces being necessary for safety concerns (i.e., prevent overcrowding) during peak times, the long distances added by the redundant spaces can be unacceptable for passengers during non-peak times. In this situation, temporary \hyperref[(P13) Shortcuts or optimizing paths]{Shortcuts or optimizing paths (P13)}, should be adopted, or make the spaces as \hyperref[(11) Adaptive redundant spaces with path regulation]{Adaptive redundant spaces with path regulation (11)}.
2) Overcrowding (or the more extreme - stampedes) is a tricky issue that is associated with many causes that are not yet fully understood by existing scientific research \citep{almeidaHumanStampedesUpdated2019}. Redundant spaces may not fully function if the crowd of people, due to crowd psychology, still concentrates in front of the destination areas (despite there being redundant spaces outside destination areas). In this situation, on-site management regulations are still needed to disperse people more evenly across the spaces.

\begin{figure}[h]
    \centering
    \includegraphics[width=0.5\columnwidth]{Paper3Imgs_pdfs/index_pictures.011.pdf}
    \caption{Redundant spaces or setting apart bottlenecks at the district level}
    \label{fig:layout_setapart}
\end{figure} \FloatBarrier
\begin{figure}[h]
    \centering
    \includegraphics[width=0.5\columnwidth]{Paper3Imgs_pdfs/index_pictures.015.pdf}
    \caption{Redundant spaces or setting apart bottlenecks at the building level}
    \label{fig:build_setapart}
\end{figure} \FloatBarrier







\subsubsection{(20) \& (21) City passages and their positioning}
\label{(20 21) City passages and their positioning}

\textbf{Context \& Problems}: It is common for railway tracks to cut city areas apart. City passages and their positioning should be designed/planned for the following purposes:
1) Facilitating passengers to move from one side of the station to the other side of the station.
2) Facilitating city users passing through the station areas that were previously blocked by railway tracks. The station areas will be more vibrant with more users passing by.
3) When stations are \hyperref[Connect with neighborhoods]{Connected with neighborhoods (8)}, the city passages across railway tracks, help to form coherent pedestrian paths at the district level, or \hyperref[(25) City axes]{City axes (25)} on an even larger scale. These paths or city axes facilitate pedestrian movement.

\textbf{Solutions \& examples}: At the district level, city passages across railway tracks should be positioned according to the city's neighborhood conditions (fig. \ref{fig:layoutpassage}). City passages may link the most dense neighborhoods/communities/city areas to maximize their usage. City passages should avoid bringing interference between transport passengers and city leisure users. If the passages are mainly for leisure users to pass by, then their positions should avoid areas where passengers intensively move. If the passages are mainly for transport passengers, their position should avoid areas where leisure users would stay or gather.

% enabling pedestrians to move freely across tracks or station buildings. This may, therefore, increase the usage of the spaces and facilities in station areas, which reduces emptiness during non-peak time. 
% (and event typology?)

If there are areas reserved for leisure events, to make events more visible or even interactive for pedestrians, the paths should be aligned close to these areas (see principle \hyperref[(6) Make events visible by aligning main paths and open spaces]{Make events visible by aligning main paths and open spaces (6)}).

At the building level, city passages can be designed with fine details, with enough lighting, and not too dark to make people feel unsafe (fig. \ref{fig:buildpassage}). Beyond city passages that link different city areas across railway tracks (city-to-city), there can be other types of passages that link station to city and station inner areas to station inner areas (e.g., the internal footbridge on platforms at the Railway Central Station) These passages also provide convenient movement and ease transfer.

Passages are commonly seen in Dutch railway stations, such as Amsterdam Central Station, Utrecht Central Station, and Rotterdam Central Station.

\textbf{Limitations}: The more passages, the more choices for pedestrians to pass through, but it takes financial cost to build passages. 
at the building level, it is often quite difficult to vertically arrange the passages with the railway platforms, the ground, and the internal passages that passengers move from platform to platform.

\begin{figure}[h]
    \centering
    \includegraphics[width=0.5\columnwidth]{Paper3Imgs_pdfs/index_pictures.023.pdf}
    \caption{Position city passages (at the district level)}
    \label{fig:layoutpassage}
\end{figure} \FloatBarrier
\begin{figure}[h]
    \centering
    \includegraphics[width=0.5\columnwidth]{Paper3Imgs_pdfs/index_pictures.013.pdf}
    \caption{The city passages (at the building level)}
    \label{fig:buildpassage}
\end{figure} \FloatBarrier






\subsubsection{(22) \& (23) Reduce barriers to ease flow}
\label{(22 23) Reduce barriers to ease flow}

\textbf{Contexts}: Barriers to flow should be reduced, especially in the following contexts: 
1) When stations suffer from emptiness, there is a need to promote the flow of people (leisure users from the city neighborhoods), letting them use the redundant station spaces for public life, leisure activities, events, and so on.
2) When \hyperref[(25) City axes]{City axes (25)} (in which the stations are located) are planned to promote the flow.
3) When stations already are \hyperref[Connect with neighborhoods]{Connected with city neighborhoods (8)}, but due to the connections' qualities being poor, only passengers driven by strong motives would pass through the connections, while leisure users are not willing to use the connections.

\textbf{Problems}: Barriers of flow are usually in the following forms with respective reasons: 
1) Stairs as barriers. Pedestrians need to make efforts to overcome level changes, and much attention is needed to walk on the steps in case of falling down.
2) Lengthy, narrow, and boring paths.
3) Tunnels, skywalks/footbridges - Pedestrians need to make efforts to overcome level changes - Pedestrians would rather walk on the ground if possible. Tunnels are usually quite dark, making pedestrians feel unsafe. Skywalks/footbridges are unpleasant to walk on if they are non-weather-protected, without shading from trees or building canopies.
4) Walking areas that are interfered with by intensive vehicle traffic - Pedestrians feel unsafe to walk through, and much attention is needed to prevent being hit by a car.
5) Strict security checking points (Common in Chinese stations) - Too much effort is needed to pass through them. 
6) Ticket gates - They make gated areas inaccessible for people without tickets.

Barriers can lead to commercial stores in station areas being shut down during non-peak times, since city leisure users, as potential customers, are not willing to overcome barriers to explore the stores.


\textbf{Solutions}: Therefore, it is necessary to Reduce barriers to ease flow, at the district level (fig. \ref{fig:beijingwestnorth}) and at the building level (fig. \ref{fig:buildflowcommercial}). The following interventions can be taken, depending on the barrier types:
1) Use ground-level paths as much as possible, to avoid level changes, and to continue the street-level qualities of the overall environment.
2) \hyperref[(9 10) Smooth level changes by landscape design]{Smooth level changes by landscape design (9) \& (10)}. (This is particularly useful with small level changes; For large level changes, using escalators is usually more appropriate).
3) For lengthy boring paths, design them like \hyperref[(7) Human-oriented spaces]{Human-oriented spaces (7)}, make them more pleasant with leisure facilities, otherwise use \hyperref[(P13) Shortcuts or optimizing paths]{Shortcuts or optimizing paths (P13)}
4) If footbridges/skywalks are adopted, create street-like qualities, such as weather-protected shading and convenience facilities.
5) Separate traffic roads from pedestrian paths, e.g., horizontally make alignment between the pedestrian flow and traffic flow to avoid overlapping, or vertically submerging traffic roads underground or elevating traffic roads.
6) Security checks and ticket gates are primarily management issues. For spatial design considerations, possibly \hyperref[(P13) Shortcuts or optimizing paths]{Shortcuts or optimizing paths (P13)} can help. 
7) Consider \hyperref[(20 21) City passages and their positioning]{City passages and their positioning (20) \& (21)} if station areas are cut apart by railway tracks.

% During non-peak times, city neighborhoods use the stations' spaces for non-transport public activities. Commercial services within the station areas will be more resilient as not only the passengers but also the city users are their potential customers.
% brings another benefit - increasing the use of commercial facilities. Given the same types and amount of commercial supply, spaces with no bottlenecks or barriers better support passengers to explore the whole variety of commercial facilities, and the commercial stores are more resilient the the local changes of users since they are patronized by passengers from the global network. 

% This reduces the redundant spaces needed, hence reducing and overcrowding as well as emptiness. 

% \textbf{Examples}:

\textbf{Limitations}: Given barriers are reduced, the flow of leisure users should not cause severe interference to passenger flow. Also, every measure has a (financial) cost, hence reducing barriers should bring enough benefits.

\begin{figure}[h]
    \centering
    \includegraphics[width=0.5\columnwidth]{Paper3Imgs_pdfs/placeholder.pdf}
    \caption{Barriers that hinder the flow at the district level}
    \label{fig:beijingwestnorth}
\end{figure} \FloatBarrier
\begin{figure}[h]
    \centering
    \includegraphics[width=0.5\columnwidth]{Paper3Imgs_pdfs/index_pictures.021.pdf}
    \caption{Reduce barriers to ease flow at the building level}
    \label{fig:buildflowcommercial}
\end{figure} \FloatBarrier







\subsubsection{(24) Stairs as stages or seats}
\label{(24) Stairs as stages or seats}

\textbf{Contexts \& Problems}: Stairs are common building components that primarily address level changes. Some stairs also serve for symbolic functions, for example, using stairs to elevate buildings (or building entrances) to make the buildings look more grand.

Typically, in station areas, stairs should be avoided as much as possible because they hinder the flow of pedestrians and may even result in people falling down. This issue is especially problematic if the flow is intensive. To address small level changes, instead of using stairs, it is better to \hyperref[(9 10) Smooth level changes by landscape design]{Smooth level changes by landscape design (9) \& (10)}. To address large level changes, instead of using stairs, it is better to use escalators. Given the same width, escalators have a larger capacity than stairs, and escalators, if not running, can still be used as stairs. 

However, in some situations, large stairs still exist in station areas: Those designed as backups of escalators for evacuations - like the grand stairs at the west of Utrecht Central Station. Or, stairs just happen to exist in the station areas - like the grand stairs of the Cologne Cathedral in front of the Cologne Central Station, and the indoor stairs in the hall of Antwerpen Central Station.

Stairs should be designed as potential stages for event performances or seats for the audience if the following two conditions are met:
1) Grand stairs, for whatever reasons, exist in railway station areas
2) Leisure events are supposed to be held in the station areas, as a measure to reduce emptiness.

% \hyperref[(6) Make events visible by aligning main paths and open spaces]{Make events visible by aligning main paths and open spaces (6)})

\textbf{Solutions}:
Grand stairs (in non-flow-intensive areas) can be used as audience seat areas, or as stages for performances during events. This takes advantage of the level change to leverage the see-and-be-seen relationship between audiences and performers in events. Such usage has been seen in various real-world cases: The grand stairs at the west side of Utrecht Central Station used during open-air exhibitions (fig. \ref{fig:utrechtstair}); The grand stairs in the hall of Antwerpen Central Station used for music festivals; And the grand stairs in front of Cologne Central Station used as stages for demonstrations.

\textbf{Limitations}: Stairs are typically unfavorable elements in the station areas, as mentioned in the above Contexts \& Problems text. If stairs are indeed used for events, the frequency of use is still questionable.




\begin{figure}[h]
    \centering
    \includegraphics[width=0.5\columnwidth]{Paper3Imgs_pdfs/index_pictures.005.pdf}
    \caption{Stairs as stages or seats}
    \label{fig:utrechtstair}
\end{figure} \FloatBarrier




\subsubsection{(25) City axes}
\label{(25) City axes}


\textbf{Context \& Problems}: If station areas suffer from emptiness, city axes, which are a common planning strategy in urban planning and urban design fields, can be planned. 


\textbf{Solutions}: City axes help to emphasize the connections, concentrate the flow of leisure users, and energize the environment.

City axes make the spatial structure of certain areas recognizable to decision-makers, thereby also promoting collaborations for different stakeholders to align the important spaces and facilities, for creating memorable experiences for users when they move in the environment.

City axes should be supported by smaller design principles to ensure pedestrians can really move within them and experience the environment. These smaller principles include: (stations) \hyperref[Connect with neighborhoods]{Connect with neighborhoods (8)}, \hyperref[(22 23) Reduce barriers to ease flow]{Reduce barriers to ease flow (22) \& (23)}, and properly design/plan \hyperref[(20 21) City passages and their positioning]{City passages and their positioning (20) \& (21)}. 

\textbf{Examples}: City axes have been successfully used in (re)development projects of some station areas, for example, the Utrecht Central Station area (fig. \ref{fig:City axes}), the Rotterdam Central Station area, and the Delft Station area.

\textbf{Limitation}: Of notice, city axes usually should be space-based (e.g., the Paris Historical Axis) instead of building-based (e.g., the axis of the Forbidden City in Beijing). Building-based city axes only work symbolically to demonstrate order, power, or spiritual values, and can hardly contribute to real space usage because people cannot move on the axes. 

\begin{figure}[h]
    \centering
    \includegraphics[width=0.5\columnwidth]{Paper3Imgs_pdfs/index_pictures.025.pdf}
    \caption{Two city axes were planned during the redevelopment of the Utrecht Central Station area}
    \label{fig:City axes}
\end{figure} \FloatBarrier

% \subsubsection{(26) Reconfigurable spaces}
% \label{(26) Reconfigurable spaces}

% adaptive redundant spaces VS path regulation
% networks of elements
% networks of point-of-interests
% networks of paths


% \subsubsection{(27) Event-supportive spaces}
% \label{(27) Event-supportive spaces}

% \subsubsection{(26) Nature solutions}
% \label{(26) Nature solutions}

% \subsubsection{(27) Compensate indoor space with more (outdoor) gray space}
% \label{(27) Compensate indoor space with more (outdoor) gray space}





\end{document}
